\documentclass[11pt,a4paper]{article}

\usepackage[margin=18mm]{geometry}
\usepackage{booktabs}
\usepackage{tabularx}
\usepackage{enumitem}
\usepackage{xcolor}
\usepackage{tcolorbox}
\usepackage{hyperref}
\usepackage{fancyhdr}
\usepackage{titlesec}
\usepackage{multirow}
\usepackage{array}
\usepackage{amssymb}

\definecolor{ink}{HTML}{0F172A}
\definecolor{accent}{HTML}{BE123C}
\definecolor{soft}{HTML}{FFF1F2}
\definecolor{muted}{HTML}{64748B}
\definecolor{ok}{HTML}{065F46}

\hypersetup{colorlinks=true,linkcolor=accent,urlcolor=accent}
\pagestyle{fancy}
\fancyhf{}
\fancyhead[L]{\small\textcolor{muted}{DA3021 · Phone Look Project}}
\fancyhead[R]{\small\textcolor{muted}{Capture checklist}}
\fancyfoot[C]{\small\textcolor{muted}{\thepage}}
\renewcommand{\headrulewidth}{0.4pt}
\renewcommand{\footrulewidth}{0.3pt}

\titleformat{\section}{\large\bfseries\color{accent}}{\thesection}{0.6em}{}
\titleformat{\subsection}{\normalsize\bfseries\color{ink}}{\thesubsection}{0.5em}{}
\setlist[itemize]{leftmargin=1.3em,itemsep=0.25em,topsep=0.3em}
\setlist[enumerate]{leftmargin=1.4em,itemsep=0.25em,topsep=0.3em}

\tcbset{
  colback=soft,
  colframe=accent,
  boxrule=0.6pt,
  arc=2pt,
  left=8pt,right=8pt,top=6pt,bottom=6pt,
  fonttitle=\bfseries
}

\begin{document}

\begin{center}
  {\LARGE\bfseries\color{ink} Capture Checklist}\\[0.35em]
  {\large\color{muted} iPhone $\leftrightarrow$ Normal\quad\&\quad Normal $\leftrightarrow$ HDR/Night}\\[0.5em]
  {\normalsize DA3021 Course Project · shoot tonight · copy into \texttt{data/raw/pairs/}}
\end{center}
\vspace{0.4em}

\begin{tcolorbox}[title=One rule for both modes]
Every scene needs \textbf{two photos of the same view}.\\
\texttt{sceneXX\_src.jpg} = source (ordinary look)\qquad
\texttt{sceneXX\_tgt.jpg} = target (look you want to copy)
\end{tcolorbox}

\section{Choose your capture mode}
You only need \textbf{one} of these. Prefer Mode~A if an iPhone (or any flagship) is available for one evening.

\begin{center}
\renewcommand{\arraystretch}{1.25}
\begin{tabular}{@{}>{\bfseries}l p{5.6cm} p{5.6cm}@{}}
\toprule
 & \textbf{Mode A · iPhone $\leftrightarrow$ Normal} & \textbf{Mode B · Normal $\leftrightarrow$ HDR/Night} \\
\midrule
\texttt{src} & Your budget / normal phone & Your phone, \textbf{normal} mode \\
\texttt{tgt} & iPhone (or Pixel / flagship) & Your phone, \textbf{Night or HDR} mode \\
Need & Friend's better phone for $\sim$1 evening & Only your phone \\
Question & ``Can my phone look like an iPhone?'' & ``Can normal shots look like Night/HDR?'' \\
Code change & None --- same filenames & None --- same filenames \\
\bottomrule
\end{tabular}
\end{center}

\section{File naming (must match the code)}
Put files in \texttt{phone\_look\_project/data/raw/pairs/}:
\begin{verbatim}
scene01_src.jpg
scene01_tgt.jpg
scene02_src.jpg
scene02_tgt.jpg
...
scene16_src.jpg
scene16_tgt.jpg
\end{verbatim}
Aim for \textbf{12--20 scenes}. More is better; fewer than 8 makes the train/test split weak.

\section{How to shoot one scene}
\begin{enumerate}
  \item Brace the phone(s): table edge, books, wall lean, or tripod.
  \item Frame a mostly \textbf{static} scene (avoid walking people in the centre).
  \item Take \texttt{src}, then within a few seconds take \texttt{tgt} from the \textbf{same spot} and framing.
  \item For Mode~A: hold the two phones stacked / side-by-side as tightly as possible.
  \item For Mode~B: switch mode only (normal $\rightarrow$ Night/HDR); do not walk between shots.
  \item Leave auto exposure / auto WB on --- those choices are part of the look.
\end{enumerate}

\section{Mode A extras (iPhone $\leftrightarrow$ normal)}
\begin{itemize}
  \item Prefer scenes a few metres away (less parallax between lenses).
  \item Match orientation (both portrait or both landscape).
  \item Keep the same zoom (both at $1\times$).
  \item If the iPhone uses Night mode automatically in the dark, that is fine --- still save it as \texttt{tgt}.
\end{itemize}

\section{Mode B extras (normal $\leftrightarrow$ HDR/Night)}
\begin{itemize}
  \item Use whichever ``best'' mode looks clearly different (Night indoor, HDR for windows).
  \item If Night needs a long exposure, keep the phone very still or the pair will not align.
  \item Do not crop differently between \texttt{src} and \texttt{tgt}.
\end{itemize}

\section{Scene checklist (tick as you go)}
Mix indoor / outdoor / faces / text / colour. Tick when both \texttt{src} and \texttt{tgt} are saved.

\vspace{0.3em}
\renewcommand{\arraystretch}{1.35}
\begin{tabularx}{\textwidth}{@{}c >{\bfseries}l X c c@{}}
\toprule
\# & Scene idea & Why it helps & src & tgt \\
\midrule
01 & Hostel desk under LED & Warm cast, noise & $\square$ & $\square$ \\
02 & Face / portrait (consent) & Skin tones & $\square$ & $\square$ \\
03 & Window / backlight & Highlights \& shadows & $\square$ & $\square$ \\
04 & Whiteboard / printed text & Sharpness / readability & $\square$ & $\square$ \\
05 & Corridor & Mixed lighting & $\square$ & $\square$ \\
06 & Food / colourful object & Saturation & $\square$ & $\square$ \\
07 & Outdoor dusk & Low light & $\square$ & $\square$ \\
08 & Plant / leaves & Fine texture & $\square$ & $\square$ \\
09 & Bookshelf & Structure \& colour & $\square$ & $\square$ \\
10 & Notice board / poster & Text + colour & $\square$ & $\square$ \\
11 & Bathroom / mirror lights & Hard light & $\square$ & $\square$ \\
12 & Stairwell & Dim indoor & $\square$ & $\square$ \\
13 & Outdoor building facade & Daylight baseline & $\square$ & $\square$ \\
14 & Desk lamp close-up & Strong local light & $\square$ & $\square$ \\
15 & Classroom / lab bench & Everyday usefulness & $\square$ & $\square$ \\
16 & Your choice & Fill gaps & $\square$ & $\square$ \\
\bottomrule
\end{tabularx}

\section{Optional bursts (extra marks)}
For 4--5 scenes, also shoot 6--8 quick frames in normal mode and save:
\begin{verbatim}
data/raw/bursts/scene01/0001.jpg
data/raw/bursts/scene01/0002.jpg
...
\end{verbatim}
Used later to show noise falling as frames are merged.

\section{After shooting --- run on the laptop}
\begin{enumerate}
  \item Copy all pairs into \texttt{data/raw/pairs/}.
  \item Delete the fake demo pairs (\texttt{scene01}--\texttt{scene08} from the smoke test) \textbf{or} use new names like \texttt{real01\_src.jpg}.
  \item Run:
\end{enumerate}
\begin{verbatim}
cd phone_look_project
python -m src.align_pairs
python -m src.fit_models
python -m src.evaluate
python -m src.make_figures
python -m src.app
\end{verbatim}

\begin{tcolorbox}[colback=white,colframe=ok,title=Quick sanity check]
After \texttt{align\_pairs}, open \texttt{data/aligned/sceneXX/}: you should see \texttt{src.png}, \texttt{tgt\_warped.png}, and \texttt{mask.png}. If alignment fails, reshoot that scene with more overlap / less motion.
\end{tcolorbox}

\section{Privacy}
\begin{itemize}
  \item Get consent before using someone else's face in the video.
  \item Blur strangers if the photo will be shown publicly.
\end{itemize}

\vspace{0.6em}
\begin{center}
\small\textcolor{muted}{Mode A and Mode B use the same code path. Only what you put in \texttt{tgt} changes.}
\end{center}

\end{document}
